


\newpage


\dspace
\subhead
1. Motifs de Tate mixte sur l'anneau des
$S$-entiers d'un corps de nombres
\endsubhead

\subhead
1.1
\endsubhead
Soit $k$ un corps. 
Nombre de r\'esultats cit\'es d\'ependant de
la r\'esolution des singularit\'es, nous
supposons $k$ de caract\'eristique $0$.
Hanamura (1995), Levine (1998) et Voevodsky
(2000) ont chacun d\'efini une cat\'egorie
triangul\'ee de motifs sur $k$, et Levine
(1998) VI 2.5.5 construit une \'equivalence
entre sa cat\'egorie triangul\'ee et celle de
Voevodsky.
Cette derni\`ere sera pour nous la plus
commode.
Nous la noterons $\DM(k)$ et noterons
$\DM(k)_{\dbQ}$ celle qui s'en d\'eduit par
tensorisation avec $\dbQ$.
On dispose dans $\DM(k)$ d'objets de Tate
$\dbZ(n)$ ($n\in\dbZ$), dont les images dans
$\DM(k)_{\dbQ}$ seront not\'ees $\dbQ(n)$.
Seule nous importera la sous-cat\'egorie
triangul\'ee $\DMT(k)_{\dbQ}$ de
$\DM(k)_{\dbQ}$ engendr\'ee par les $\dbQ(n)$:
celle des ``extensions it\'er\'ees'' de
$\dbQ(n)[m]$.
Rappelons que $E$ est {\it extension} de $B$
par $A$ s'il existe un triangle distingu\'e
$A\to E\to B\to A[1]$.

On dispose dans $\DM(k)$ d'un produit
tensoriel associatif et commutatif \`a unit\'e
$\otimes$, compatible \`a la structure
triangul\'ee.
L'automorphisme de sym\'etrie de
$\dbZ(1)\otimes\dbZ(1)$ est l'identit\'e
(Voevodsky (2000) 2.1.5), $M\mapsto
M(1):=M\otimes \dbZ(1)$ est une \'equivalence
(loc. cit. 4.1.3), et $\dbZ(n)$ est
$\dbZ(1)^{\otimes n}$ ($n\in\dbZ$).
Le produit tensoriel $\otimes$ est rigide
(loc. cit. 4.3.7): existence pour tout objet
$M$ d'un dual $M^*$ muni de $\ev\colon\,
M^*\otimes M\to \dbZ(0)$ et $\delta\colon\,
\dbZ(0)\to M\otimes M^*$ tels que les compos\'es
$M\to M\otimes M^*\otimes M\to M$ et $M^*\to
M^*\otimes M\otimes M^*\to M^*$ soient l'identit\'e.

Rappelons, d'apr\`es Levine (1993), comment,
lorsque la conjecture d'annulation de
Beilinson Soul\'e est v\'erifi\'ee
pour $k$, on peut extraire de $\DMT(k)_{\dbQ}$
une cat\'egorie tannakienne $\MT(k)$ qui
m\'erite le nom de cat\'egorie des motifs de
Tate mixte sur $k$.

Ecrivons $\Hom^j(M,N)$ pour $\Hom(M,N[j])$.
Les $\Hom^*(\dbZ(a),\dbZ(b))$ ne d\'ependant
que de $i=b-a$.
Ils sont donn\'es par les groupes de Chow
sup\'erieurs convenablement num\'erot\'es de
$\Spec(k)$ (loc. cit. 4.2.9 pour
$X=\Spec(k)$).
Pour $i< 0$ ils sont nuls. 
Pour $i=0$,
$$
\align
&\matrix
\Hom^j(\dbZ(0),\dbZ(0))
   &=\dbZ &\qquad \text{ si $j=0$},\\
   &\hfill 0 &\qquad \text{ sinon}.\hfill
\endmatrix\tag1.1.1\\
\intertext{Pour $i=1$,}
&\matrix
\Hom^j(\dbZ(0),\dbZ(1))&=k^* 
   &\qquad\text{ si $j=1$,}\\
   &\,\,\,\,0 &\qquad \text{ sinon}.\hfill
\endmatrix\tag1.1.2
\endalign
$$

Apr\`es tensorisation avec $\dbQ$, ces groupes
sont encore donn\'ees par les sous-espaces
propres des op\'erations d'Adams dans les
groupes de $K$-th\'eorie de $k$ (Levine
(1998) II 3.6.6; r\'ef\'erence originale:
Bloch (1986) compl\'et\'e par Bloch (1994);
voir aussi Levine (1994)).
La conjecture d'annulation de
Beilinson-Soul\'e est que
$$
\Hom^j(\dbQ(0),\dbQ(i))=0\quad
\text{pour}\quad i>0\quad\text{et}\quad
j\Le 0.
\tag1.1.3
$$
Lorsque tel est le cas, Beilinson Bernstein
Deligne (1982) 1.3.14, appliqu\'e \`a la
sous-cat\'egorie pleine de $\DMT(k)_{\dbQ}$
d'objets les $\dbQ(n)$, montre que les
extensions it\'er\'ees de $\dbQ(n)$
($n\in\dbZ$) forment une cat\'egorie
ab\'elienne, coeur d'une $t$-structure sur
$\DMT(k)_{\dbQ}$.
On l'appelle la cat\'egorie
$\MT(k)$ des motifs de Tate mixte sur $k$.
Elle est stable par produit tensoriel et
passage au dual.
Ces foncteurs \'etant exacts, cela r\'esulte
de ce que $\dbQ(n)\otimes\dbQ(m)=\dbQ(n+m)$ et
que $\dbQ(n)^*=\dbQ(-n)$.
On obtient ainsi sur $\MT(k)$ un produit
tensoriel exact en chaque variable,
associatif, commutatif, \`a unit\'e et rigide.

De ce que $\MT(k)$ est le coeur d'une
$t$-structure sur $\DMT(k)_{\dbQ}$ r\'esulte
que pour $A$ et $B$ dans $\MT(k)$, on a
$$
\align
\Ext^1(A,B) &=
\Hom^1(A,B)\qquad\text{et}\tag1.1.4\\
\Ext^2(A,B) &\hookrightarrow \Hom^2(A,B).\tag1.1.5
\endalign
$$
En particulier,
$$
\Ext^1(\dbQ(n),\dbQ(m))=0\quad
\text{pour}\quad m\Le n.\tag1.1.6
$$
Pour $A$ dans $\MT(k)$, il r\'esulte de
(1.1.6), par induction sur une description de
$A$ comme extension it\'er\'ee, que $A$ admet
une unique filtration ``par le poids''
$W$, finie, croissante et index\'ee par les
entiers pairs, telle que
$$
\Gr_{-2n}^W(M):= W_{-2n}(M)/W_{-2(n+1)}(M)
$$
soit une somme de copies de $\dbQ(n)$.
Posons
$$
\omega_n(M):=\Hom(\dbQ(n),\Gr_{-2n}^W(M)).
$$
On a donc 
$$
\Gr_{-2n}^W(M)=\dbQ(n)\otimes_\dbQ
\omega_n(M).
$$
La filtration $W$ est fonctorielle, exacte et
compatible au produit tensoriel.
Parce que l'automorphisme de symm\'etrie de
$\dbQ(1)\otimes\dbQ(1)$ est l'identit\'e, on
dispose d'isomorphismes 
$\dbQ(n)\otimes\dbQ(m)\simover\dbQ(n+m)$
compatibles \`a l'associativit\'e et \`a la
commutativit\'e de $\otimes$.
Le foncteur exact $M\mapsto
\omega(M):=\oplus\omega_n(M)$ est donc
un $\otimes$-foncteur: c'est un foncteur fibre
et la cat\'egorie $MT(k)$ est tannakienne.

\subhead
1.2
\endsubhead
Pour $M$ un motif de Tate mixte, le
sous-quotient $W_{-2n}(M)/W_{-2(n+2)}(M)$ de
$M$ est une extension $M_{n,n+1}$ de
$\Gr_{-2n}^W(M)=\dbQ(n)\otimes\omega_n(M)$ par
$\Gr_{-2(n+1)}^W(M)=\dbQ(n+1)\otimes\omega_{n+1}(M)$.
La classe de cette extension est un \'el\'ement
$$
e_n\in \Ext^1(\dbQ(n),\dbQ(n+1))\otimes\Hom
(\omega_n(M),\omega_{n+1}(M)).
$$
La tensorisation avec $\dbQ(n)$ est un
isomorphisme
$$
\Ext^1(\dbQ(0),\dbQ(1))\simover
\Ext^1(\dbQ(n),\dbQ(n+1))\tag1.2.1
$$
et les $e_n$ d\'efinissent
$$
e_M\colon\,\omega_*(M)\to\Ext^1(\dbQ(0),
\dbQ(1))\otimes\omega_*(M).
\tag1.2.2
$$

Pour $x\colon\,\dbQ(n)\to\Gr_{-2n}^W(M)$ dans
$\omega_n(M)$, l'image inverse par $x$ de
l'extension $M_{n,n+1}$ est une extension de
$\dbQ(n)$ par $\Gr_{-2(n+1)}^W(M)=\dbQ(n+1)
\otimes\omega_{n+1}(M)$, et $e_M(x)$ est la
classe de cette extension dans
$$
\Ext^1(\dbQ(0),\dbQ(1))\otimes
\omega_{n+1}(M)=\Ext^1(\dbQ(n),
\dbQ(n+1))\otimes\omega_{n+1}(M).
$$

On regardera $e_M$ comme une coaction, au sens
des coalg\`ebres de Lie.
Plus pr\'ecis\'ement, comme une coaction, sur
$\omega(M):=\oplus \omega_*(M)$, de la
coalg\`ebre de Lie librement coengendr\'ee par
$\Ext^1(\dbQ(0),\dbQ(1))$.

\proclaim{Proposition 1.3}
La coaction $e_M$ est fonctorielle en $M$ et
compatible au produit tensoriel.
\endproclaim

\demo{Preuve}
La v\'erification de la fonctorialit\'e en $M$
est laiss\'ee au lecteur.
Pour $x\in\omega_m(M)$ et $y\in\omega_n(N)$, on
doit montrer que
$$
e_{M\otimes N}(x\otimes y)=e_M(x)\otimes y
+x\otimes e_N(y).
\tag1.3.1
$$
Par fonctorialit\'e, on peut pour le
v\'erifier remplacer $M$ par $W_{-2m}(M)$, une
extension de $\dbQ(m)\otimes\omega_m(M)$ par
$W_{-2(m+1)}(M)$, puis par l'image inverse de
cette extension par
$x\colon\,\dbQ(m)\to\dbQ(m)\otimes\omega_m(M)$.
De m\^{e}me pour $N$, $n$ et $y$.
On se ram\`ene ainsi \`a supposer que
$M=W_{-2m}(M)$, $\omega_m(M)=\dbQ$, $x=1$, et
de m\^{e}me pour $N$, $n$ et $y$.
On a alors
$$
\align
\omega_{m+n+1}(M\otimes N)\simover
  &\omega_{m+n+1}(M\otimes\text{(quotient
$\dbQ(n)$ de $N$)})\\
\oplus\,\, &\omega_{m+n+1}\text{((quotient
$\dbQ(m)$ de $M$)}\otimes N).
\endalign
$$
et, par fonctorialit\'e en $M$ et $N$ de
$e_{M\otimes N}$, il suffit de v\'erifier
(1.3.1) pour $M$ remplac\'e par son quotient
$\dbQ(m)$, ainsi que pour $N$ remplac\'e par
son quotient $\dbQ(n)$.
Ces cas se r\'eduisent \`a la d\'efinition de
l'isomorphisme (1.2.1).
\enddemo

Si $E$ est une extension de $\dbQ(0)$ par
$\dbQ(1)$, il r\'esulte de 1.3 que la classe
de l'extension $E^*(1)$ de $\dbQ(0)$ par
$\dbQ(1)$ est l'oppos\'ee de la classe de $E$.

\definition{Definition 1.4}
{\it Pour $\Gamma$ un sous-espace vectoriel de
$\Ext^1(\dbQ(0),\dbQ(1))$, 
$MT(k)_\Gamma$ est la cat\'egorie des motifs de
Tate mixte tels que la coaction (1.2.2) se
factorise par une coaction de $\Gamma$}.
\enddefinition

Pour que $M$ soit dans $\MT(k)_\Gamma$, il
faut et il suffit que pour tout sous-quotient
$E$ de $M$ qui est une extension d'un
$\dbQ(n)$ par $\dbQ(n+1)$, la classe de $E$
dans 
$$
\Ext^1(\dbQ(0),\dbQ(1)\simover \Ext^1
(\dbQ(n),\dbQ(n+1))
$$ 
soit dans $\Gamma$. 
D'apr\`es 1.3, cette sous-cat\'egorie de
$MT(k)$, qui contient l'objet unit\'e $\dbQ(0)$,
est stable par sous-quotients, produits
tensoriels et duaux.
C'est une sous-cat\'egorie tannakienne de
$MT(k)$.

\subhead
1.5
\endsubhead
Quel que soit $k$ (de caract\'erisque $0$), on
dispose d'un foncteur ``r\'ealisation'' de
$\DM(k)$ dans la cat\'egorie triangul\'ee
``syst\`emes de r\'ealisations'' de Huber (1995).
L'\'etat de la litt\'erature \`a ce sujet
n'est pas satisfaisant, sp\'ecialement en ce qui
concerne la r\'ealisation de Hodge.
Dans Huber (2000), il est suppos\'e \`a tort
dans la preuve de 2.1.4, p.~775 que les
hyperrecouvrements propres de $X$ forment une
cat\'egorie filtrante.
Cette erreur est corrig\'ee dans un erratum
dat\'e d'octobre 2002.
Dans Levine (1998) V. 2.3.7 le foncteur Dec
de d\'ecalage d'une filtration par le poids
est appliqu\'e \`a un complexe de faisceaux,
alors qu'il devrait \^{e}tre appliqu\'e
apr\`es passage \`a $R\Gamma$.

Nous esquissons ci-dessous une fa\c{c}on
possible de proc\'eder, adapt\'ee \`a la
d\'efinition de Voevodsky
de $\DM(k)$, que nous commen\c{c}ons par
rappeler.
Le foncteur obtenu sera contravariant
(cohomologique).
Dans la suite de l'article, c'est sa variante
covariante (homologique), d\'eduite par
passage au dual, qui nous sera utile.

Soit $\SmCor(k)$ la cat\'egorie d'objets les
sch\'emas s\'epar\'es lisses sur $k$, un
morphisme de $X$ dans $Y$ \'etant une
correspondance finie de $X$ vers $Y$:
$\Hom(X,Y)$ est le groupe ab\'elien librement
engendr\'e par les sous-sch\'emas r\'eduits
irr\'eductibles ferm\'es $Z$ de $X\times Y$,
finis sur $X$ et dominant une composante
connexe de $X$.
La cat\'egorie $\SmCor(k)$ est additive
(avec $\oplus=\Coprod$).
La cat\'egorie triangul\'ee $\DM(k)$ se d\'eduit
de la cat\'egorie triangul\'ee de complexes
born\'es $K^b(\SmCor(k))$ par

\medskip\noindent
{\sans (a)} localisation:
on divise par la sous-cat\'egorie \'epaisse
engendr\'ee par les $[X\times A^1]\to[X]$
(invariance 
d'homotopie) et les $[U\CAP V]\to[U]\oplus[V]
\to[X]$ pour $X$ la r\'eunion des ouverts $U$
et $V$ (Mayer-Vietoris); 

\smallskip\noindent
{\sans (b)}\enspace 
adjonction de facteurs directs images
d'endomorphismes idempotents;

\smallskip\noindent
{\sans (c)}\enspace
inversion formelle de $\dbQ(1)$.

\medskip
Il r\'esulte de Voevodsky (2000) 
3.2.6 et de (a), Mayer
Vietoris, qu'on obtiendrait une cat\'egorie
\'equivalente si dans $\SmCor(k)$ on se
limitait aux sch\'emas lisses et
quasi-projectifs sur $k$.
Pour d\'efinir le foncteur r\'ealisation,
l'essentiel est de d\'efinir la r\'ealisation
d'un complexe born\'e $X^*$ d'objets de
$\SmCor(k)$, et par la remarque ci-dessus, on
peut se ramener au cas o\`u les $X^i$ sont
quasi-projectifs.

\proclaim{Lemme 1.5.1}
Soit $\Gamma\colon\, X\to Y$ dans $\SmCor(k)$.
Si $\Ybar$ est une compactification projective
et lisse de $Y$, il existe une
compactification projective et lisse $\Xbar$
de $X$ telle que $\Gamma$ se prolonge en
$\Gammabar\colon\,\Xbar\to\Ybar$.
quand elle existe, l'extension $\Gammabar$ 
est unique.
\endproclaim

\demo{Preuve}
Il suffit de traiter du cas o\`u $X$ est
connexe et o\`u $\Gamma$ est un sous-sch\'ema
ferm\'e $Z$ de $X\times Y$, r\'eduit,
irr\'eductible, fini et dominant sur $X$.
Si $d$ est le degr\'e de $Z$ sur $X$, $Z$ d\'efinit
un morphisme de sch\'emas $z$ de $X$ dans
$\Sym^d(Y)$.
D'apr\`es Hironaka, il existe une
compactification projective et lisse $\Xbar$
de $X$ telle que $z$ se prolonge en
$\zbar\colon\, \Xbar\to\Sym^d(\Ybar)$.
Pour un tel $\Xbar$, l'adh\'erence $\Zbar$ de
$Z$ dans $\Xbar\times\Ybar$ est encore finie
sur $\Xbar$.
Elle fournit le prolongement cherch\'e.
L'unicit\'e est claire.
Noter que, par Hironaka encore, on peut
supposer que $X$ est le compl\'ement dans
$\Xbar$ d'un diviseur \`a croisements normaux.

Pour $X^*$ un complexe born\'e d'objets de
$\SmCor(k)$, on d\'eduit de 1.5.1 l'existence
d'un syst\`eme de compactifications $\Xbar^n$
des $X^n$ tel que \ (a) $\Xbar^n$ est
projectif et lisse, et $X^n$ est le
compl\'ement dans $\Xbar^n$ d'un diviseur \`a
croisements normaux $D^n$; \ (b) la
diff\'erentielle $d$ se prolonge de $X^*$ \`a
$\Xbar^*$.
Par unicit\'e du prolongement, le prolongement
v\'erifie encore $d^2=0$.
Les syst\`emes de compactifications $\Xbar$
forment une cat\'egorie filtrante.
Pour $f\colon\, X^*\to Y^*$ un morphisme de
complexes, et pour $\Ybar^*$ une
compactification de $Y^*$ v\'erifiant (a),
(b), il existe une compactification $\Xbar^*$
de $X^*$ v\'erifiant (a), (b) et telle que $f$
se prolonge en un morphisme de complexes de
$\Xbar^*$ dans $\Ybar^*$.
\enddemo

D\'efinissons la r\'ealisation de de Rham de
$X^*$ -- ingr\'edient essentiel de la
r\'ealisation de Hodge relative \`a un
plongement complexe de $k$ dans $\dbC$.
L'id\'ee est de choisir une compactification
$\Xbar^*$ de $X^*$ comme ci-dessus et de
prendre sur chaque $\Xbar^n$ le complexe de de
Rham \`a p\^{o}les logarithmiques
$\Omega_{\Xbar^n}^*(\log\,D^n)$.
Ce complexe est bifiltr\'e: par le filtration
de Hodge $F^p=\sigma_{\Ge p}$, et par la
filtration par le poids $W_*$, qui compte le
nombre de facteurs $\frac{dz_i}{z_i}$ pour
$z_i$ une \'equation locale d'une composante
lisse de $D^n$.
Il s'agit ensuite d'obtenir des complexes
bifiltr\'es $(K^{(n)},W,F)$ repr\'esentant les
$R\Gamma(\Xbar^n,\Omega^*(\log\,D^n))$ et
d'utiliser la diff\'erentielle de $\Xbar^*$
pour en faire un complexe double bifiltr\'e,
avec $K^{(n)}$ comme colonne de premier indice
$-n$.
La difficult\'e sera de d\'eduire de $d^2=0$
pour $\Xbar^*$ que $d'{^2}=0$ identiquement.
Ceci fait, il reste \`a

\medskip\noindent
{\sans (i)}\enspace
remplacer la filtration $W$ de $K^{(n)}$ par
son translat\'e $W'$ d\'efini par
$W'{^m}=W^{m-n}$;

\smallskip\noindent
{\sans (ii)}\enspace
passer au complexe simple associ\'e; on notera
que son $\Gr^W$ est la somme des\break
$\Gr^W(K^{(n)})[n]$;

\smallskip\noindent
{\sans (iii)}\enspace
remplacer la filtration $W$ par sa d\'ecal\'ee
$\Dec\,W$.

\medskip
Nous travaillerons avec la topologie \'etale.
Celle de Nisnevitch ferait aussi bien
l'affaire, mais non celle de Zariski.
Pour un faisceau quasi-coh\'erent, la
cohomologie est la m\^{e}me.
L'hypercohomologie d'un complexe de de Rham
aussi.
Gain: pour $p\colon\, Z\to X$ un morphisme
fini, le foncteur $p_*$ est exact.
Ceci est faux pour Zariski.
Soit $\Gamma\colon\, X\to Y$ dans $\SmCor(k)$:
$\Gamma=\sum n_i Z_i$ et $Z_i\Subset X\times
Y$ est fini sur $X$.
Soient des compactifications $X$ et $Y$ en $\Xbar$ et
$\Ybar$, avec \`a l'infini un diviseur \`a
croisement normaux, et supposons que $\Gamma$
se prolonge en $\Gammabar$, i.e. que chaque
adh\'erence $\Zbar_i$ de $Z_i$ dans
$\Xbar\times\Ybar$ soit finie sur $\Xbar$.
Soit enfin $\vert\Gammabar\vert$ le support
$\cup \Zbar_i$ de $\Gammabar$.
La correspondance $\Gammabar$ d\'efinit alors
un morphisme de complexes de faisceaux sur
$\vert\Gamma\vert$
$$
[\Gammabar]\colon\, \pr_2^*\Omega^*
(\log)\to\pr_1^*\Omega^*(\log).
$$
Ici aussi que la topologie soit \'etale ou de
Nisnevitch est essentiel.
Prendre garde que $\pr_2^*$ est une image
inverse au sens des faisceaux, et non pas au
sens des faisceaux coh\'erents (pour laquelle
l'image inverse de la diff\'erentielle de de
Rham $d$ ne serait pas d\'efinie).
Ce morphisme s'\'etend aux r\'esolutions
flasques canoniques $\scrC^*$ et il suffit de
repr\'esenter
$R\Gamma(\Xbar^n,\Omega^*(\log))$ par
$\Gamma(\Xbar^n,\scrC^*\Omega^*(\log))$.

\subhead
1.6
\endsubhead
En caract\'eristique z\'ero, les seuls
corps pour lesquels la conjecture d'annulation
de Beilinson-Soul\'e soit d\'emontr\'ee sont
les corps de nombres ($=$ extensions finies de
$\dbQ$), les corps de fonctions rationelles
d'une courbe de genre $0$ sur un corps de
nombres, et les limites inductives de tels
corps.
Le cas qui nous importe est celui des corps de
nombres.
Pour ceux-ci, $K_n(k)\otimes\dbQ$ est nul pour
$n$ pair $>0$ et on a
$$
\align
\Ext^1(\dbQ(0),\dbQ(n))
&=K_{2n-1}(k)\otimes\dbQ,\quad\text{et}\\
\Ext^2(\dbQ(0),\dbQ(n)) &=0,
\endalign
$$
car $\Hom^2(\dbQ(0),\dbQ(n))=0$ dans
$\DM(k)_{\dbQ}$.

Soit $k$ un corps de nombres.
Le foncteur ``r\'ealisation'' sur $\DM(k)$
induit sur $\MT(k)$ un $\otimes$-foncteur \`a
valeurs dans les syst\`emes de r\'ealisations
au sens de Deligne (1989) (aspect cristallin
exclus) ou Jannsen (1990).

Rappelons que
$\Ext^1(\dbQ(0),\dbQ(1))=k^*\otimes\dbQ$.
Si $S$ est un ensemble de places finies de
$k$, et que $\scrO_S$ est l'ensemble des
$S$-entiers de $k$ (int\'egralit\'e en dehors
de $S$), nous d\'efinissons artificiellement la
cat\'egorie des motifs de Tate mixte sur
$\scrO_S$ comme \'etant $\MT(k)_\Gamma$
pour $\Gamma=\scrO_S^*\otimes\dbQ\Subset
k^*\otimes\dbQ^*$.
Cas particulier: pour $S$ le compl\'ement
d'une place finie $v$, $\scrO_S$
est le localis\'e $\scrO_{(v)}$ en $v$ de
l'anneau des entiers $\scrO$ de $k$, et on
parle de motifs de Tate mixte {\it non ramifi\'es}
en $v$.
Un motif de Tate mixte sur $\scrO_S$ est donc
un motif de Tate mixte sur $k$ non ramifi\'e
en chaque place finie $v\notin S$.

\proclaim{Proposition 1.7}
Soit $\ell$ un nombre premier distinct de la
caract\'eristique r\'esiduelle de $v$.
Pour qu'un motif de Tate mixte $M$ sur $k$ soit
non ramifi\'e en $v$, il faut et il suffit que
sa r\'ealisation $\ell$-adique $M_\ell$
soit non ramifi\'ee en $v$.
\endproclaim

\demo{Preuve}
La r\'ealisation $\ell$-adique est un
$\dbQ_\ell$-espace vectoriel d\'ependant
fonctoriellement du choix d'une cl\^{o}ture
alg\'ebrique $\kbar$ de $k$.
La fonctorialit\'e en $\kbar$ en fait une
repr\'esentation continue de $\Gal(\kbar/k)$.
Nous utiliserons que 

\medskip\noindent
{\sans (a)}
la r\'ealisation l-adique de $\dbQ(1)$
est $\dbQ_\ell(1)
:=\dbZ_\ell(1)\otimes_{\dbZ_\ell}\dbQ_\ell$,
avec $\dbZ_\ell(1)=\lim\,\mu_{\ell^n}(\kbar)$;

\medskip\noindent
{\sans (b)}
Pour $n\not=m$, puisque $\Hom(\dbQ(n),\dbQ(m))=0$,
une extension de $\dbQ(n)$ par $\dbQ(m)$ est
d\'etermin\'ee \`a isomorphisme unique pr\`es
par sa classe dans $\Ext^1(\dbQ(n),\dbQ(m))$.
Pour $x\in k^*\Subset\Ext^1(\dbQ(0),\dbQ(1))$,
notons $K(x)$ l'extension de $\dbQ(0)$
par $\dbQ(1)$, dite de Kummer, de classe $x$.
Sa r\'ealisation l-adique est d\'eduite,
par tensorisation avec $\dbQ_\ell$, de
l'extension  de Kummer $K_\ell(x)$ de
$\dbZ_\ell$ par $\dbZ_\ell(1)$.
La donn\'ee d'une extension
$$
0\to\dbZ_\ell(1)\to E\mapright{f}
\dbZ_\ell\to 0
$$
de  $\dbZ_\ell$ par $\dbZ_\ell(1)$ \'equivaut
\`a celle du $\dbZ_\ell(1)$-torseur
$f^{-1}(1)$, et $K_\ell(x)$ correspond au
$\dbZ_\ell(1)$ torseur limite projective des
$\mu_{\ell^n}(\kbar)$-torseurs des racines
$(\ell^n)^{\iemes}$ de $x$ dans $\kbar$.

\medskip
La repr\'esentation $\ell$-adique $K_\ell(x)$
est non ramifi\'ee en $v$ si et seulement si
$x$ est une unit\'e en $v$, i.e. est dans
$\scrO_{(v)}^*$.
Pour $M$ de Tate mixte, il r\'esulte donc de
(b) que $M$ est dans $MT(\scrO_v)$ si et
seulement si les repr\'esentations
$\ell$-adiques
$W_{-2n}(M_\ell)/W_{-2(n+2)}(M_\ell)$ sont non
ramifi\'ees.
Il reste \`a v\'erifier le lemme suivant.

\proclaim{Lemme 1.8}
Soient $K_v$ le corps des fractions de
l'hens\'elis\'e $\scrO_{(v)}^h$ de $\scrO_v$
et $H$ une repr\'e-\break
sentation $\ell$-adique de
$\Gal(\Kbar_v/K_v)$, munie d'une filtration
finie croissante index\'ee par les entiers
pairs $W$, telle que $\Gr_{-2n}^W(H)$ soit
somme de copies de $\dbQ_\ell(n)$.
Pour que la repr\'esentation $H$ soit non
ramifi\'ee, il suffit que les
repr\'esentations $W_{-2n}/W_{-2(n+2)}$ le
soient.
\endproclaim

\demo{Preuve}
Soit $I_v\Subset \Gal(\Kbar_v/K_v)$ le groupe
d'inertie.
On sait que son plus grand pro-$\ell$-quotient
est canoniquement isomorphe \`a
$\dbZ_\ell(1)$.
Soit $t_\ell\colon\, I_v\to\dbZ_\ell(1)$ le
morphisme de passage au quotient.
Il est $\Gal(\Kbar_v/K_v)$-\'equivariant.
\enddemo

Supposons que $I_v$ agisse trivialement sur
les $W_{-2n}/W_{-2(n+2)}$ et prouvons par
r\'e-\break
currence sur $r\Ge 2$ qu'il agit
trivialement sur les $W_{-2n}/W_{-2(n+r)}$.
Pour $r=2$, c'est l'hypo-\break
th\`ese.
Pour $r>2$, l'hypoth\`ese  de r\'ecurrence
assure que l'action de $I_v$ est triviale sur
$W_{-2n}/W_{-2(n+r-1)}$ et sur
$W_{-2(n+1)}/W_{-2(n+r)}$.
L'action de $\sigma\in I_v$ est donc de la
forme $1+\nu(\sigma)$, avec pour $\nu(\sigma)$
un compos\'e
$$
W_{-2n}/W_{-2(n+r)}\to
\Gr_{-2n}^W\longmapright{35}{n(\sigma)}
\Gr_{-2(n+r-1)}^W\to W_{-2n}/W_{-2(n+r)},
$$
et $n(\sigma_1\sigma_2)=n(\sigma_1)+n(\sigma_2)$.
L'application $\sigma\mapsto n(\sigma)$
se factorise donc par $\dbZ_\ell(1)$, et est
de la forme
$$
n(\sigma)=n \cdot t_\ell(\sigma),\,\,
\text{\rm pour }n\colon\,\Gr_{-2n}^W(1)\to
\Gr_{-2(n+r-1)}^W.
$$
Le morphisme $n$ est
$\Gal(\Kbar_v/K_v)$-equivariant.
Puisque $r>2$,
$\Hom_{\Gal(\Kbar_v/K_v)}(\dbQ_\ell(n+1~),$\break
$\dbQ_\ell(n+r-1))=0$ et $n$ est nul, ce qui
conclut la r\'ecurrence.
\enddemo

\proclaim{Proposition 1.9}
Soit $\Gamma$ un sous-espace vectoriel de
$\Ext^1(\dbQ(0),\dbQ(1))$.

\medskip\noindent
{\sans (i)}
Pour $r\Ge 2$, $\Ext^1(\dbQ(n),\dbQ(n+r))$ est
le m\^{e}me dans $MT(k)$ et dans
$MT(k)_\Gamma$.

\noindent
{\sans (ii)}
Les $\Ext^2$ de Yoneda sont nuls dans
$MT(k)_\Gamma$.
\endproclaim

\demo{Preuve}\quad (i) Si $r\Ge 2$, une
extension de $\dbQ(n)$ par $\dbQ(n+r)$ est en
effet automatiquement dans $MT(k)_\Gamma$.

\noindent
(ii)
Il suffit de v\'erifier la nullit\'e des
$\Ext^2(\dbQ(n),\dbQ(m))$.
Pour toute classe $c$ dans cet $\Ext^2$, il
existe dans $MT(k)_\Gamma$ une extension
$$
0\to F\to E_1\to\dbQ(n)\to 0,
\tag1.8.1
$$
de classe not\'ee $c_1$, et une extension
$E_2$ de $F$ par $\dbQ(m)$, de classe $c_2$,
telle que $c$ soit le produit de Yoneda
$c_2c_1$.
Ce produit est l'image de $c_2$ par le cobord
$$
\partial\colon\, \Ext^1(F,\dbQ(m))\to
\Ext^2(\dbQ(n),\dbQ(m))
$$
d\'efini par (1.8.1).
C'est aussi l'obstruction \`a l'existence dans
$MT(k)_\Gamma$ de\break
$X=X_0\Supset X_1\Supset
X_2\Supset X_3=0$  de quotients successifs
$X_i/X_{i+1}$ ($i=0,1,2$) munis\break
d'isomorphismes avec $\dbQ(n)$, $F$ et
$\dbQ(m)$, tel que la classe de l'extension
$X/X_2$ de $\dbQ(n)$ par $F$ soit $c_1$, et
que la classe de l'extension $X_1$ de $F$ par
$\dbQ(m)$ soit $c_2$.

Relevons $\dbQ(n)$ dans  $\Gr_{-2n}^W(E_1)$ et
soit $E'_1$ l'image inverse dans
$W_{-2n}(E_1)$ de ce rel\`evement.
C'est une extension de $\dbQ(n)$ par $F':=
F\CAP E'_1=W_{-2(n+1)}(F)$.
Si $E'_2$ est l'image inverse de $F'$ dans
$E_2$, $c$ est encore le produit de Yoneda
$c'_2c'_1$ de la classe $c'_1$ de $E'_1$ par
la classe $c'_2$ de $E'_2$.
On a gagn\'e que $F'=W_{-2(n+1)}(F')$.
Un argument dual permet ensuite de se ramener
au cas o\`u $F=W_{-2(n+1)}(F)$ et
$W_{-2m}(F)=0$.
Le $\Ext^2$ ne peut donc \^{e}tre non nul que
si $m\Ge n+2$.
Puisque le $\Ext^2$ est nul dans $MT(k)$, il
existe dans $MT(k)$ une extension it\'er\'ee
$X=X_0\Supset X_1\Supset X_2\Supset X_3=0$ du
type consid\'er\'e plus haut, le sous-quotient
$W_{-2a}/W_{-2(a+2)}$ de $X$ coincide avec le
m\^{e}me sous-quotient de $X_0/X_3$ si $a+2\le
m$, et avec le m\^{e}me sous-quotient de $X_1$
si $a>n$.
Si, comme on peut le supposer, $m\Ge n+2$, on
est pour tout a dans un de ces deux cas, $X$
est donc dans $MT(k)_\Gamma$, et $c=0$.
\enddemo





